\documentclass[10pt,a4paper]{amsart}
\usepackage[margin=19mm]{geometry}
\usepackage{amsmath,amssymb,mathtools}
\usepackage[scr=rsfs,cal=euler]{mathalfa}
\usepackage{xfrac}
\usepackage[unicode,hidelinks,bookmarks=false]{hyperref}
\newcommand{\ie}{i.e.\ }
\newcommand{\cf}{cf.\ }
\newcommand{\resp}{respectively\ }
\newcommand{\Subsection}{\subsection}
\providecommand{\nobreakdash}{\nobreak\hbox{-}}
\setlength{\emergencystretch}{2em}
\multlinegap0pt
\usepackage{bm}
\usepackage{bbm}
\usepackage{dsfont}
\usepackage{xspace}
\usepackage{cancel}
\usepackage{mathtools}
\usepackage{amsthm}
\makeatletter
\@ifundefined{definition}{\newtheorem{definition}{Definition}}{}
\@ifundefined{lemma}{\newtheorem{lemma}{Lemma}}{}
\@ifundefined{proposition}{\newtheorem{proposition}{Proposition}}{}
\@ifundefined{theorem}{\newtheorem{theorem}{Theorem}}{}
\makeatother
\newcommand{\vol}{{\rm vol}}
\title[Note on edge expansion and modularity in preferential attach]{Note on edge expansion and modularity in preferential attachment graphs}
\author{Colin McDiarmid, Katarzyna Rybarczyk, Fiona Skerman,, Ma{\l}gorzata Sulkowska}
\date{Source: arXiv:2601.05953v1 (CC BY 4.0)}
\begin{document}
\maketitle

\noindent\emph{Source and licence.} Note on edge expansion and modularity in preferential attachment graphs. Version arXiv:2601.05953v1 (CC BY 4.0), distributed under the Creative Commons Attribution 4.0 International licence, as recorded in the arXiv licence metadata for this exact version and shown on the abstract page https://arxiv.org/abs/2601.05953v1. Full rights notice: licenses/arxiv-2601.05953-CC-BY-4.0.txt.\par\smallskip

\noindent\textit{Editorial scope.} Counted unit: the Lemma whose source label is \texttt{lem:dull} in the article (source0.tex lines 463--466, source label \texttt{lem:dull}), delivered together with the article's own complete original proof of it (source0.tex lines 468--511, one \texttt{proof} environment of 44 lines). No mathematics is written, repaired or re-derived: statement and proof are the article's own text line for line. Required context retained uncounted: thm:general\_upper\_bound\_mod\_2 (lines 239--246); defn.neg\_mod (lines 434--439); prop.bound\_by\_worst\_set\_vertex (lines 441--446). Every definition, display and label of those lines is retained unchanged. Omitted from this package and not redistributed: 1--238, 247--433, 440--440, 447--462, 467--467, 512--753. Nothing inside a retained excerpt is omitted or altered: every retained body is delivered exactly as the article's lines are written, so no omission receipt of any kind accompanies this package. Citations: the retained excerpts carry no citation command at all, so no bibliography entry of the article is transcribed and no reference is redistributed. Reference adaptation: none was needed and none was made; every reference inside the delivered unit resolves inside it, so no unresolved reference or citation is produced. Printed slips kept literal: no printed word, symbol, hypothesis or number of the counted statement or of its proof was altered. Layout and class: the article's own class is \texttt{llncs} with packages including fontenc, amsmath, amssymb, graphicx, enumitem, comment, url; the wrapper is the lane's pinned \texttt{amsart} editorial wrapper at 10pt on A4, which restates the theorem-like environments the retained text needs and adds the article's own macro definitions that the retained text needs (\texttt{\textbackslash vol}), with extra wrapper packages (bm, bbm, dsfont, xspace, cancel, mathtools). The delivered numbering is the unit's own. Assets: no retained span contains a figure environment, a graphics-inclusion command, a PDF-page inclusion, a table or a TikZ picture; no image file is copied into this package, the package declares no asset dependency and no image file is redistributed with it. Licence: version arXiv:2601.05953v1 of this article is distributed under the Creative Commons Attribution 4.0 International licence, as recorded in the arXiv licence metadata for this exact version and shown on the abstract page https://arxiv.org/abs/2601.05953v1. A full rights notice is delivered at licenses/arxiv-2601.05953-CC-BY-4.0.txt. Characters: every retained excerpt is pure ASCII, so the delivered engine, XeLaTeX with its default Latin Modern text font, reports no missing character.
\begin{theorem}\label{thm:general_upper_bound_mod_2}
Let $0<u\leq 1/2$, and $G=(V,E)$ be a graph with $u$-expansion $\alpha_u(G)$, minimum degree $h \geq 1$, and such that $e(S) \leq h |S|$ for all $S \subseteq V$. 
Define   
$\delta_u:= \min\{\alpha_u(G)/h,1\}$. Then
	\[
	q^*(G) \leq 1- \inf_{0 < u \leq 1/2} \left(\frac{\delta_u}{2+\delta_u} + \frac{u}{2}\right).
	\]
\end{theorem}

\begin{definition}\label{defn.neg_mod}
	Let $G=(V,E)$ be a graph with at least one edge, and let $S \subseteq V$.	Define the negative relative modularity of $S$ to be 
	\[
	q_{vr}^{-}(S) = \frac{\vol(G)}{\vol(S)}\left(\frac{ e(S,\bar{S})}{2e(G)}+\frac{\vol(S)^2}{\vol(G)^2}\right). 
	\]
\end{definition}

\begin{proposition} \label{prop.bound_by_worst_set_vertex}
	Let $G=(V,E)$ be a graph with at least one edge and let $\mathcal{A}$ be any vertex partition of $G$. Then
    \[
		q_{\mathcal{A}}(G) \leq 1 - \min_{S \in \mathcal{A}} q_{vr}^{-}(S).
	\]
\end{proposition}

\begin{lemma}\label{lem:dull}
	Let $0 <  u \leq 1/2$, and $0 \leq \delta \leq 1$. Then
	\[ \frac{\delta}{2+\delta} + \frac{u}{2} \leq \frac{\delta u}{2(1-u) +\delta u} + \frac{1-u}{2}.\]
\end{lemma}

\begin{proof}[of Theorem~\ref{thm:general_upper_bound_mod_2}]
	We write $G=(V,E)$ and $n=|V|$ and we will use Proposition \ref{prop.bound_by_worst_set_vertex}. We first show that for any $S \subseteq V$ such that $|S| =un$, where $0<u\leq 1/2$ we get
	\[
	q_{vr}^-(S) \geq \delta_u/(2+\delta_u)+u/2
	\]
	(see Defintion~\ref{defn.neg_mod}) and consider larger sets later. By the assumption of the theorem we have
	\[
		\vol(S) = 2e(S) + e(S,\bar{S}) \leq 2 h|S| + e(S,\bar{S})
	\]
	and by the definition of $u$-expansion
	\[
		e(S,\bar{S}) \geq \alpha_u(G)|S| = \delta_u h |S|.
	\]
	Therefore
	\begin{equation}\label{eq.boundedgeloss}
		\frac{e(S, \bar{S})}{\vol(S)} \geq \frac{e(S, \bar{S})}{2 h|S| + e(S,\bar{S})} 
		\geq \frac{\delta_u h|S|}{2h|S|+ \delta_u h|S|} = \frac{\delta_u}{2+\delta_u}.
	\end{equation}
    Thus for $S$ with $|S| =  un \leq n/2$ by~\eqref{eq.boundedgeloss} and since $\vol(S)\geq h|S|$,
    
    \[ q_{vr}^{-}(S) = \frac{e(S, \bar{S})}{\vol(S) } + \frac{\vol(S)}{\vol(G)} \geq \frac{\delta_u}{2+\delta_u} + \frac{h|S|}{2hn} =  \frac{\delta_u}{2+\delta_u} + \frac{u}{2}, \]
    hence we get 
   
    \begin{equation}\label{eq.usual_sets}\min_{\substack{S\subseteq V\\|S|\leq n/2}} q_{vr}^{-}(S)\geq \inf_{0 < u \leq 1/2} \Big( \frac{\delta_u}{2+\delta_u}+\frac{u}{2} \Big).
	\end{equation}
	%
	Now consider $S \subseteq V$ such that $|S| = \bar{u}n$ and $1/2<\bar{u}\leq 1$. Note that $|\bar{S}|=(1-\bar{u})n \leq n/2$ and so, by the $u$-expansion definition, $e(S,\bar{S}) \geq \delta_{1-\bar{u}}h|\bar{S}|$. This time, similarly to~\eqref{eq.boundedgeloss},
	%
	\[
		\frac{e(S, \bar{S})}{\vol(S)} 
		%
		\geq \frac{e(S, \bar{S})}{2 h|S| + e(S,\bar{S})}
		\geq \frac{\delta_{1-u} h|\bar{S}|}{2h|S|+ \delta_{1-\bar{u}} h|\bar{S}|} = \frac{\delta_{1-\bar{u}} (1-\bar{u})}{2\bar{u}+\delta_{1-\bar{u}}(1-\bar{u})}
	\]
	and 
	\[ q_{vr}^{-}(S) = \frac{e(S, \bar{S})}{\vol(S) } + \frac{\vol(S)}{\vol(G)} \geq \frac{\delta_{1-\bar{u}}(1-\bar{u})}{2\bar{u}+\delta_{1-\bar{u}}(1-\bar{u})} + \frac{h|S|}{2hn} =  \frac{\delta_{1-\bar{u}}(1-\bar{u})}{2\bar{u}+\delta_{1-\bar{u}}(1-\bar{u})} + \frac{\bar{u}}{2}. \]
	Writing $s=1-\bar{u}$, for $S \subseteq V$ such that $|S|=\bar{u}n > n/2$ we get
	\[ q_{vr}^{-}(S)  \geq \frac{\delta_{s}s}{2(1-s)+\delta_{s}s} + \frac{1-s}{2} \geq \frac{\delta_s}{2+\delta_s}+ \frac{s}{2} \]
	where the second inequality followed by Lemma~\ref{lem:dull}. Thus the case when $|S| > n/2$ is covered by our existing bounds in~\eqref{eq.usual_sets} for $|S|\leq n/2$. We get
	\[
		\min_{S \subseteq V} q_{vr}^{-}(S)\geq \inf_{0 < u \leq 1/2} \Big( \frac{\delta_u}{2+\delta_u}+\frac{u}{2} \Big)
	\]
	and the conclusion follows by Proposition~\ref{prop.bound_by_worst_set_vertex}.
\end{proof}

\end{document}
